\section{Additional Templates for the CoWM Study}
\label{app:budget_extensions}
These panels reserve measurements required by specific resource or deployment claims. They contain no measured results and do not imply that conditional experiments have been completed.

\paragraph{Sharing and resource cost.}
Table~\ref{tab:resources} compares the resources of separate and shared backbones. Both total resident and active parameters are counted, including the visual encoder. Training steps, data exposure, precision, and hardware accompany per-update time and total GPU-hours. A same-width separate-DiT system has greater total capacity; an equal-capacity comparison is a distinct optional control.
\begin{table*}[t]
\centering\scriptsize
\resizebox{\textwidth}{!}{\begin{tabular}{lrrlrrrr}
\toprule
Configuration & Total M & Active M & Updates & Step ms & Train GiB & GPU-hours & Decision ms \\
\midrule
Separate-Recorded & \PaperResult{draft-resource-separate-total} & \PaperResult{draft-resource-separate-active} & \PaperResult{draft-resource-separate-updates} & \PaperResult{draft-resource-separate-step-ms} & \PaperResult{draft-resource-separate-memory} & \PaperResult{draft-resource-separate-hours} & \PaperResult{draft-resource-separate-p50} \\
Shared-Recorded & \PaperResult{draft-resource-recorded-total} & \PaperResult{draft-resource-recorded-active} & \PaperResult{draft-resource-recorded-updates} & \PaperResult{draft-resource-recorded-step-ms} & \PaperResult{draft-resource-recorded-memory} & \PaperResult{draft-resource-recorded-hours} & \PaperResult{draft-resource-recorded-p50} \\
Shared-Pred-Full & \PaperResult{draft-resource-full-total} & \PaperResult{draft-resource-full-active} & \PaperResult{draft-resource-full-updates} & \PaperResult{draft-resource-full-step-ms} & \PaperResult{draft-resource-full-memory} & \PaperResult{draft-resource-full-hours} & \PaperResult{draft-resource-full-p50} \\
\bottomrule
\end{tabular}
}
\caption{Resource measurements for completed sharing contrasts. Decision latency uses matched $N/S/K$ and the same hardware. Sharing weights alone does not remove the separate A/B computations.}
\label{tab:resources}
\end{table*}

\paragraph{Readouts.}
Table~\ref{tab:physical_probe} reports linear and MLP readouts of observed latents. The main Cube panel covers all eight physical quantities; three auxiliary task panels use task-native state fields. Targets and units are fixed from the environment schema, and both models use identical episode-disjoint probe rows. These measurements test information recoverable by a readout, not the accuracy of future-state prediction. Applying the frozen observed-latent readouts to predicted futures remains a separate extension.
\begin{table*}[t]
\centering\small
\resizebox{\textwidth}{!}{\begin{tabular}{llrrrrrrrr}
\toprule
Quantity & Target / units & \multicolumn{4}{c}{LeWM} & \multicolumn{4}{c}{CoWM} \\
 & & \multicolumn{2}{c}{Linear} & \multicolumn{2}{c}{MLP} & \multicolumn{2}{c}{Linear} & \multicolumn{2}{c}{MLP} \\
 & & MSE$\downarrow$ & $r\uparrow$ & MSE$\downarrow$ & $r\uparrow$ & MSE$\downarrow$ & $r\uparrow$ & MSE$\downarrow$ & $r\uparrow$ \\
\midrule
\multicolumn{10}{l}{\textbf{Cube (main)}} \\
Block position & \texttt{\detokenize{privileged_block_0_pos}} (meters) & $0.0080$ & $0.9960$ & $0.0018 \pm 0.0001$ & $0.9991 \pm 0.0001$ & $0.0102$ & $0.9949$ & $0.0010 \pm 0.0001$ & $0.9995 \pm 0.0000$ \\
Block quaternion ($w/x/y/z$) & \texttt{\detokenize{privileged_block_0_quat}} (unitless) & $0.7581$ & $0.0335$ & $0.8299 \pm 0.0081$ & $0.0315 \pm 0.0139$ & $0.7065$ & $0.1829$ & $0.7126 \pm 0.0008$ & $0.2013 \pm 0.0040$ \\
Block yaw & \texttt{\detokenize{privileged_block_0_yaw}} (radians) & $0.9698$ & $0.0231$ & $1.1134 \pm 0.0129$ & $0.0228 \pm 0.0122$ & $0.9649$ & $-0.0863$ & $0.9838 \pm 0.0015$ & $0.0048 \pm 0.0109$ \\
End-effector position & \texttt{\detokenize{proprio_effector_pos}} (meters) & $0.0201$ & $0.9899$ & $0.0027 \pm 0.0000$ & $0.9987 \pm 0.0000$ & $0.0020$ & $0.9990$ & $0.0006 \pm 0.0001$ & $0.9997 \pm 0.0000$ \\
End-effector yaw & \texttt{\detokenize{proprio_effector_yaw}} (radians) & $1.1062$ & $-0.0011$ & $1.3113 \pm 0.0405$ & $0.0076 \pm 0.0168$ & $1.0744$ & $-0.0068$ & $1.1105 \pm 0.0387$ & $0.0013 \pm 0.0645$ \\
Gripper opening & \texttt{\detokenize{proprio_gripper_opening}} (unitless normalized fraction) & $0.0763$ & $0.9611$ & $0.0139 \pm 0.0006$ & $0.9932 \pm 0.0002$ & $0.0114$ & $0.9943$ & $0.0085 \pm 0.0001$ & $0.9957 \pm 0.0001$ \\
Joint position & \texttt{\detokenize{proprio_joint_pos}} (radians) & $0.2211$ & $0.6737$ & $0.2371 \pm 0.0110$ & $0.6773 \pm 0.0047$ & $0.1975$ & $0.6972$ & $0.2019 \pm 0.0039$ & $0.6934 \pm 0.0081$ \\
Joint velocity & \texttt{\detokenize{proprio_joint_vel}} (radians/second) & $1.0122$ & $0.1354$ & $1.0189 \pm 0.0077$ & $0.1625 \pm 0.0093$ & $0.6624$ & $0.5238$ & $0.4479 \pm 0.0017$ & $0.7151 \pm 0.0017$ \\
\textbf{Overall} & Equal weight over 8 attributes & $0.5215$ & $0.4764$ & $0.5661 \pm 0.0071$ & $0.4866 \pm 0.0039$ & $0.4536$ & $0.5374$ & $0.4333 \pm 0.0052$ & $0.5764 \pm 0.0100$ \\
\addlinespace
\multicolumn{10}{l}{\textbf{TwoRoom (auxiliary)}} \\
Agent position & \texttt{\detokenize{pos_agent}} (environment coordinate units (pixels)) & $0.0074$ & $0.9964$ & $0.0001 \pm 0.0000$ & $0.9999 \pm 0.0000$ & $0.0008$ & $0.9996$ & $0.0001 \pm 0.0000$ & $1.0000 \pm 0.0000$ \\
\textbf{Overall} & Equal weight over 1 attributes & $0.0074$ & $0.9964$ & $0.0001 \pm 0.0000$ & $0.9999 \pm 0.0000$ & $0.0008$ & $0.9996$ & $0.0001 \pm 0.0000$ & $1.0000 \pm 0.0000$ \\
\addlinespace
\multicolumn{10}{l}{\textbf{PushT (auxiliary)}} \\
Agent location & \texttt{\detokenize{state[0:2]}} (pixels) & $0.0371$ & $0.9805$ & $0.0023 \pm 0.0004$ & $0.9988 \pm 0.0002$ & $0.0079$ & $0.9959$ & $0.0008 \pm 0.0000$ & $0.9996 \pm 0.0000$ \\
Block location & \texttt{\detokenize{state[2:4]}} (pixels) & $0.0223$ & $0.9877$ & $0.0022 \pm 0.0005$ & $0.9988 \pm 0.0003$ & $0.0250$ & $0.9863$ & $0.0044 \pm 0.0003$ & $0.9976 \pm 0.0001$ \\
Block yaw & \texttt{\detokenize{state[4:5]}} (radians) & $0.2122$ & $0.8988$ & $0.0714 \pm 0.0086$ & $0.9672 \pm 0.0041$ & $0.2683$ & $0.8702$ & $0.0778 \pm 0.0025$ & $0.9638 \pm 0.0012$ \\
\textbf{Overall} & Equal weight over 3 attributes & $0.0905$ & $0.9557$ & $0.0253 \pm 0.0028$ & $0.9883 \pm 0.0013$ & $0.1004$ & $0.9508$ & $0.0277 \pm 0.0009$ & $0.9870 \pm 0.0004$ \\
\addlinespace
\multicolumn{10}{l}{\textbf{Reacher (auxiliary)}} \\
Joint position & \texttt{\detokenize{qpos}} (radians) & $0.3903$ & $0.7107$ & $0.3976 \pm 0.0077$ & $0.6978 \pm 0.0092$ & $0.3882$ & $0.7116$ & $0.3949 \pm 0.0092$ & $0.7018 \pm 0.0122$ \\
Joint velocity & \texttt{\detokenize{qvel}} (radians/second) & $0.9970$ & $0.0156$ & $0.9963 \pm 0.0002$ & $0.0192 \pm 0.0031$ & $0.9963$ & $0.0212$ & $0.9958 \pm 0.0003$ & $0.0278 \pm 0.0045$ \\
\textbf{Overall} & Equal weight over 2 attributes & $0.6937$ & $0.3631$ & $0.6969 \pm 0.0039$ & $0.3585 \pm 0.0059$ & $0.6922$ & $0.3664$ & $0.6954 \pm 0.0045$ & $0.3648 \pm 0.0039$ \\
\addlinespace
\bottomrule
\end{tabular}
}
\caption{Observed-latent physical readouts. Entries report standardized-coordinate MSE and coordinate-mean Pearson $r$; MLP values are mean $\pm$ sample standard deviation over three readout seeds, while Linear is deterministic. Overall is an equal-weight mean over attributes. Cube is the main panel; TwoRoom, Push-T, and Reacher are auxiliary. Probe test episodes are disjoint from readout training and validation; this does not assert that the pretrained world models never saw those episodes.}
\label{tab:physical_probe}
\end{table*}

\paragraph{Optional real-robot evaluation.}
A real-robot extension would compare success and latency with matched initialization, control frequency, action interfaces, and trial counts (Table~\ref{tab:robot}). Policy time and observation-to-command time are measured separately. Completion time specifies how failed or truncated trials are treated. This is a design template, not evidence for real-world performance.
\begin{table*}[t]
\centering\scriptsize
\resizebox{\textwidth}{!}{\begin{tabular}{llrrrrrr}
\toprule
Method & Task & Trials & Successes & SR / CI & Policy ms & System p95 ms & Completion s \\
\midrule
LeWM + CEM & \PaperResult{draft-robot-lewm-task} & \PaperResult{draft-robot-lewm-trials} & \PaperResult{draft-robot-lewm-successes} & \PaperResult{draft-robot-lewm-sr} & \PaperResult{draft-robot-lewm-policy} & \PaperResult{draft-robot-lewm-system} & \PaperResult{draft-robot-lewm-completion} \\
CoWM selection & \PaperResult{draft-robot-selection-task} & \PaperResult{draft-robot-selection-trials} & \PaperResult{draft-robot-selection-successes} & \PaperResult{draft-robot-selection-sr} & \PaperResult{draft-robot-selection-policy} & \PaperResult{draft-robot-selection-system} & \PaperResult{draft-robot-selection-completion} \\
CoWM refinement & \PaperResult{draft-robot-refinement-task} & \PaperResult{draft-robot-refinement-trials} & \PaperResult{draft-robot-refinement-successes} & \PaperResult{draft-robot-refinement-sr} & \PaperResult{draft-robot-refinement-policy} & \PaperResult{draft-robot-refinement-system} & \PaperResult{draft-robot-refinement-completion} \\
\bottomrule
\end{tabular}
}
\caption{Optional real-robot template; no trials are asserted. SR/CI denotes success rate and its uncertainty.}
\label{tab:robot}
\end{table*}

\paragraph{Search replacement criterion.}
The historical prespecified criterion requires a macro one-sided 95\% success-difference bound above $-2$ percentage points and each declared task above $-4$ points, together with at least 20\% lower p50 decision latency. A new comparison must declare its task set, action semantics, and reference budget before applying the criterion. The historical reference did not pass the complete task-level gate. Lower latency alone therefore does not establish search replacement.

\paragraph{Historical inference snapshot.}
Table~\ref{tab:reference_inference_snapshot} preserves the previous manuscript's measured fixed-checkpoint panel alongside its unfilled matched-study fields. These protocols remain separate from the new main tables.
\begin{table*}[t]
\centering\small
\begin{tabular}{lrrrrr}
\toprule
Method & Cube & Push-T & Reacher & TwoRoom & Macro \\
\midrule
\multicolumn{6}{l}{Measured fixed-checkpoint reference (one training seed)} \\
P0 & 100.0 & 95.5 & 80.0 & 99.0 & 93.6 \\
Random-64 & 100.0 & 97.5 & 82.0 & 98.0 & 94.4 \\
P3 & 100.0 & 96.5 & 89.5 & 99.5 & 96.4 \\
P0+PO & 100.0 & 95.5 & 85.0 & 99.5 & 95.0 \\
P0+GF & 100.0 & 95.5 & 83.0 & 100.0 & 94.6 \\
P2 CEM (300$\times$30) & 83.5 & 94.5 & 91.0 & 100.0 & 92.2 \\
P2-small (64$\times$5) & 88.0 & 95.0 & 83.5 & 100.0 & 91.6 \\
\midrule
\multicolumn{6}{l}{Matched-training endpoints (counts reported separately)} \\
P0 & --- & \PaperResult{infer-p0-pusht} & \PaperResult{infer-p0-reacher} & --- & \PaperResult{infer-p0-macro} \\
Random-64 & --- & \PaperResult{infer-random-pusht} & \PaperResult{infer-random-reacher} & --- & \PaperResult{infer-random-macro} \\
P3 & --- & \PaperResult{infer-p3-pusht} & \PaperResult{infer-p3-reacher} & --- & \PaperResult{infer-p3-macro} \\
P0+PO & --- & \PaperResult{infer-po-pusht} & \PaperResult{infer-po-reacher} & --- & \PaperResult{infer-po-macro} \\
P0+GF (conditional) & --- & \PaperResult{infer-gf-pusht} & \PaperResult{infer-gf-reacher} & --- & \PaperResult{infer-gf-macro} \\
\bottomrule
\end{tabular}

\caption{Preserved historical inference snapshot. The measured panel uses one training seed and its original evaluation cohort; pending rows retain their original result slots.}
\label{tab:reference_inference_snapshot}
\end{table*}
